\documentclass{article}

\usepackage[preprint]{neurips_2026}

\usepackage[utf8]{inputenc}
\usepackage[T1]{fontenc}
\usepackage{hyperref}
\usepackage{url}
\usepackage{booktabs}
\usepackage{amsfonts}
\usepackage{amsmath}
\usepackage{nicefrac}
\usepackage{microtype}
\usepackage{xcolor}

\title{Measurement-Theoretic Active Elicitation for Predicting AI
Evaluation Outcomes}

\author{%
  Dinesh Katupputhur Ramprasath \\
  Stanford University \\
  \texttt{din1993@stanford.edu} \\
}

\begin{document}

\maketitle

\begin{abstract}
Predicting whether an AI system answers a benchmark item correctly, under a
strict budget on how many ground-truth labels may be purchased, is a
measurement problem in disguise: the correctness probability is a
Rasch/IRT interaction between a subject's latent \emph{ability} and an item's
\emph{difficulty}, and the choice of which labels to buy is exactly the
efficient, Fisher-information-guided elicitation problem of adaptive testing.
We propose to recast the NeurIPS~2026 Predictive AI Evaluation Competition in
this measurement-theoretic frame. Our contribution is twofold: (i) a Bayesian
IRT predictor whose posterior over subject ability is updated as labels arrive,
yielding calibrated uncertainty where our current black-box tabular predictor
has none; and (ii) a Fisher-information / expected-Brier-reduction acquisition
policy that replaces an ad-hoc uncertainty heuristic and directly targets the
competition's area-under-the-learning-curve metric. We will evaluate against our
own 18 leaderboard submissions (best Brier~ALC $=0.162$) using an offline replay
harness and live leaderboard checks.
\end{abstract}

\section{Overview and Motivation}

\paragraph{The task.} The NeurIPS~2026 Predictive AI Evaluation Competition asks
a predictor to output $P(\text{subject answers item correctly})$ for
\mbox{(subject, item)} pairs, where a \emph{subject} is an AI system under fixed
evaluation settings and an \emph{item} is a benchmark task. The competition is a
\emph{streaming test-time adaptation} problem: for each subject--benchmark pair
the predictor may purchase ground-truth labels under a nested budget schedule
$\{0,1,3,7,15,31\}$, and is scored by \textbf{Brier~ALC} --- the (weighted) area
under the Brier learning curve across those six budgets, with weights
$(0.1,0.2,0.2,0.2,0.2,0.1)$. Lower is better; a constant $0.5$ predictor scores
$0.25$.

\paragraph{Two coupled preference-learning problems.} In the language of
CS329H, the task decomposes into (1) \emph{estimation}: predicting a latent
``will this subject get this item right,'' which is a factor model /
Rasch interaction $\sigma(\theta_s - \beta_i)$ between subject ability $\theta_s$
and item difficulty $\beta_i$ (Lectures~4--6); and (2) \emph{elicitation}:
choosing which labels to buy under a budget to drive down Brier fastest, which
is efficient Fisher-information-guided elicitation / adaptive testing
(Lecture~7).

\paragraph{The gap in our prior submissions.} Our team has submitted 18 versions
to this competition, with a current best Brier~ALC of $0.162$ (a tabular
in-context model with text embeddings). Two structural weaknesses remain, both
squarely in the scope of this course:
\begin{itemize}
  \item \textbf{The predictor is a black box.} We fit an in-context tabular
  model on the acquired labels and blend toward a hand-coded prior with a linear
  trust weight $\min(n/10, 1)$. There is no latent-trait model and no posterior
  over ability, so we cannot reason about \emph{how much} a label will help
  before buying it.
  \item \textbf{Acquisition is a one-line heuristic.} Our acquisition rule buys a
  label whenever $|\hat p - 0.5| \le 0.20$. This is a crude proxy for item
  informativeness; it ignores how much a label transfers to the frozen
  evaluation pool and ignores the budget schedule --- yet acquisition is what
  the ALC metric rewards most.
\end{itemize}

\paragraph{Research questions.}
\textbf{(RQ1, predictor)} Does a Bayesian IRT predictor with a posterior over
ability updated as labels arrive match or beat the black-box tabular predictor
on Brier, while providing calibrated uncertainty?
\textbf{(RQ2, elicitation)} Does a Fisher-information / expected-Brier-reduction
acquisition policy beat the $|\hat p - 0.5|$ heuristic at equal budgets, as
measured by Brier~ALC?
\textbf{(RQ3, cold start)} How well can the budget-0 prior be predicted from
subject/item features alone, and how much of the ALC gain comes from a better
prior versus better acquisition?

\paragraph{Why this matters.} Predicting evaluation outcomes under a label budget
is the core problem of \emph{efficient AI evaluation}: running fewer, better-chosen
evaluations to rank systems reliably \citep{truong2025amortized}. Framing it with
measurement theory yields an acquisition policy that is interpretable and
budget-aware rather than a heuristic, and connects a practical competition to the
psychometric foundations of the course. This is a well-scoped extension of
existing author infrastructure for a solo project.

\section{Connection to Course Topics}

The project uses factor and reward models (Lecture~4) to frame correctness as an
ability$\times$difficulty interaction; maximum-likelihood parameter learning from
binary outcomes (Lecture~5); Bayesian estimation with a Gaussian prior on ability
(Lecture~6), whose MAP shrinkage at low budgets is a principled replacement for
our ad-hoc trust blend; and, centrally, efficient preference
elicitation (Lecture~7): the result that the most informative item for estimating
ability is the one on which the subject has a $\approx\!50/50$ chance
(Fisher information $p(1-p)$ is maximized at $p=0.5$) is exactly what our
uncertainty heuristic approximates. We replace it with the principled quantity and
extend it from ``information about this item'' to ``expected Brier reduction over
all unlabeled evaluation items for this subject.''

\section{Literature Review}

\paragraph{Item response theory and adaptive testing.} IRT models a binary
response as $\sigma(\theta - \beta)$ and underlies computerized adaptive testing,
where items are chosen to maximize Fisher information at the current ability
estimate \citep{lord1980, embretson2000}. \citet{martinezplumed2019} brought IRT
to machine learning, estimating latent classifier ability and instance difficulty
and showing that accuracy conflates the two; \citet{zhou2026} applied IRT to LLM
benchmarks and found that hard-instance performance is more predictive of real
capability than aggregate accuracy. These establish the estimation half of our
task but do not address \emph{budgeted} label acquisition.

\paragraph{Efficient evaluation and active learning.}
\citet{truong2025amortized} predict evaluation outcomes to avoid running full
benchmarks, motivating the competition itself. Closest to the \emph{elicitation}
half of our task, \citet{polo2024tinybenchmarks} and \citet{vivek2024anchor} use
IRT and representativeness to pick a small, \emph{fixed} anchor subset that
estimates benchmark accuracy from far fewer items; our setting differs in that the
subset is chosen \emph{adaptively per subject} under a nested budget and scored
online by Brier~ALC rather than fixed offline. Active learning
\citep{settles2009} supplies the acquisition design space --- uncertainty
sampling (which our heuristic approximates) versus expected-error-reduction, the
principled policy we adopt and specialize to the Brier~ALC objective.

\paragraph{Tabular foundation models.} Our predictor backbone is an in-context
tabular transformer \citep{hollmann2023tabpfn, ye2025tabicl}, which gives strong
low-shot prediction but no latent-trait structure or calibrated posterior.

\paragraph{Gap.} Competition entries, ours included, treat the predictor as a
black box and acquisition as an uncertainty heuristic. No entry models the task as
IRT with a Bayesian posterior over ability, and none uses a Fisher-information /
expected-Brier-reduction acquisition policy tuned to the nested budget schedule.
That is this project's contribution.

\section{Proposed Methods and Analysis}

\paragraph{Predictor: Bayesian IRT with a learned prior.} We model
$P(\text{correct} \mid s, i) = \sigma(\theta_s - \beta_i)$, optionally the 2PL
form $\sigma(\alpha_i(\theta_s - \beta_i))$. At budget~0 we regress $\beta_i$ and
an ability prior from features (our existing scalar features plus text embeddings
of item content: provider tier, reasoning effort, recency), replacing the
hand-coded prior. At budget~$\ge 1$ we treat $\theta_s$ as unknown with a Gaussian
prior $\theta_s \sim \mathcal{N}(\mu,\sigma^2)$ and a Bernoulli-logit likelihood
over acquired labels, and compute a MAP / Laplace posterior per call. We estimate
$\theta_s$ \emph{per (subject, benchmark) pair} --- the unit of analysis
throughout --- not globally, so the 31-label cap buys information about one pair's
ability at a time. The posterior mean gives the prediction and the posterior
variance gives the
uncertainty that drives acquisition --- a principled replacement for the
$\min(n/10,1)$ blend. We retain the tabular model as an ensemble/residual member
and compare IRT-only, tabular-only, and blended variants.

\paragraph{Acquisition: Fisher-optimal and budget-aware.} We compare three
policies: (a) the current $|\hat p - 0.5|\le 0.20$ heuristic; (b) a
\emph{Fisher} policy that acquires the candidate maximizing $I(i)=p_i(1-p_i)$
under the current posterior, with an adaptive threshold calibrated to
\texttt{labels\_remaining}/\texttt{items\_remaining}; and (c) a novel
\emph{expected-Brier-reduction} policy that buys the label maximizing the expected
reduction in Brier over the frozen evaluation pool of the same subject--benchmark
pair. Because evaluation items share $\theta_s$, a label that sharpens the
$\theta_s$ posterior improves every evaluation prediction, directly targeting ALC
rather than single-item uncertainty. Because $\beta_i$ is itself regressed with
error at budget~0, the Fisher term $I(i)=p_i(1-p_i)$ treats $\beta_i$ as known ---
a standard CAT approximation we note and could later correct with a plug-in
difficulty variance. All policies respect the nested-prefix
schedule and the cap of 31 labels per pair.

\paragraph{Evaluation.} We build an offline replay harness from the sample data
and our assembled 24k-row offline pool: hold out each pair's evaluation pool,
replay the acquisition stream, and compute Brier at each budget and the weighted
ALC, so we can iterate without waiting on the leaderboard. To prevent leakage, the
budget-0 difficulty and prior regression is fit only on pairs disjoint from the
held-out evaluation pools it is later scored on. \textbf{Metrics:}
Brier at each budget; Brier~ALC (primary); expected calibration error (ECE); label
efficiency (ALC vs.\ labels used); and an ablation separating prior quality from
acquisition quality (RQ3). We also report IRT goodness-of-fit (infit/outfit item
statistics and a one- vs.\ multi-dimensional comparison) to test whether a single
$\theta_s$ per pair is adequate, since the pool-level acquisition argument relies
on it; and we report bootstrap 95\% confidence intervals for ALC over held-out
pairs, so that any improvement over v7 is tested for significance rather than read
off a single scalar. We validate the best offline configuration on the live
leaderboard against our v7 baseline ($0.162$).

\paragraph{Baselines.} Our prior submissions: v1 (tabular, $0.228$), v3 (tabular
$+$ prior $+$ heuristic acquisition, $0.230$), and v7 (tabular $+$ embeddings,
$\mathbf{0.162}$, current best). Success means beating v7's ALC with the IRT $+$
Fisher-acquisition stack, or matching it with a far more interpretable, calibrated
model.

\paragraph{Expected results.} Because the current heuristic (buying when
$|\hat p - 0.5|\le 0.20$, i.e.\ $p\in[0.3,0.7]$) is already a thresholded
Fisher-information rule, we expect policy~(b) to track it closely; the pool-level
expected-Brier-reduction policy~(c) is the genuine source of any gain. We expect
(1) the Fisher / expected-Brier-reduction policies to beat the uncertainty
heuristic most at small budgets (1, 3), where item choice matters and where ALC
weight concentrates; (2) the Bayesian posterior to improve low-budget calibration
(ECE) over the linear trust blend; and (3) the ablation to show acquisition
dominating gains at budgets 1--7 and the prior dominating at budget~0.
In keeping with a pre-analysis plan, we commit in advance to the following null
interpretations: if policy~(b) does not beat heuristic~(a) at any budget, we will
conclude that single-item information is already captured by the heuristic and
attribute any ALC gain to the pool-level policy~(c) and the calibrated prior; and
if the Bayesian posterior does not reduce ECE, we will report the measurement
framing as a negative calibration result rather than select metrics post hoc. We
state plainly that the scientific contribution is the acquisition policy and the
calibration analysis even if the IRT stack only ties v7 on raw ALC.

\section{Timeline, Responsibilities, and Risks}

\begin{table}[h]
  \centering
  \caption{Week-by-week plan (Pre-Analysis Plan due Nov~4; final manuscript and
  code due Dec~3).}
  \label{tab:timeline}
  \begin{tabular}{lll}
    \toprule
    Week & Dates & Milestone \\
    \midrule
    1 & Nov 5--9   & Offline replay harness reproducing leaderboard ALC; validate vs.\ v7 \\
    2 & Nov 10--16 & Bayesian IRT predictor (MAP/Laplace posterior); compare vs.\ tabular \\
    3 & Nov 17--23 & Fisher and expected-Brier-reduction policies; ablation vs.\ heuristic \\
    4 & Nov 24--30 & Cold-start prior regression (RQ3); leaderboard submission; begin writeup \\
    5 & Dec 1--3   & Final manuscript, figures, code cleanup; submit Dec~3 \\
    \bottomrule
  \end{tabular}
\end{table}

\textbf{Responsibilities.} As a solo project, I will build the predictor (IRT
posterior, feature/prior regression), the acquisition policies, and the offline
replay harness, and run the leaderboard submissions.

\textbf{Risks and mitigation.} (i)~\emph{Offline replay may not match the
leaderboard} --- calibrate against v7's known score before trusting offline
numbers and keep at least one leaderboard check per policy. (ii)~\emph{IRT may
underperform the tabular model on raw Brier} --- use IRT for the posterior/uncertainty
that drives acquisition while keeping the tabular model as the point predictor.
(iii)~\emph{Expected-Brier-reduction may be too slow for the 8-hour runtime} ---
use a closed-form Laplace approximation and fall back to Fisher $p(1-p)$.
(iv)~\emph{Originality} --- the measurement-theoretic predictor and acquisition
policy are the novel contribution; competition infrastructure is cited as the
author's own prior work, and scope will be confirmed with course staff.

\bibliographystyle{plainnat}
\bibliography{references}

\end{document}
